\documentclass[10pt,letterpaper]{article}
\usepackage{cornell-notes}

\title{Clause 04: Conformance}
\author{}
\date{}

\setDocTitle{Clause 04: Conformance}
\setDocSubtitle{{C 2024}}
\setDocOwner{{C 2024 Cornell Notes}}
\setDocDate{{\today}}
\setCornellCollection{{C 2024 Cornell Notes}}
\setCornellUnitType{{Clause}}
\setCornellUnitNumber{04}
\setCornellUnitTitle{Conformance}
\hypersetup{pdftitle={Clause 04: Conformance},pdfsubject={C 2024 study notes},pdfauthor={}}

\begin{document}
\maketitle



\begin{CornellOverviewBox}{Clause Orientation}
This clause assigns obligations to programs and implementations. It explains the force of mandatory wording, the relationship between constraints and undefined behavior, the strict-conformance portability target, the hosted and freestanding implementation categories, permitted extensions, and required implementation documentation.
\end{CornellOverviewBox}

\section{Learning Objectives}
\begin{itemize}
  \item Interpret mandatory and prohibitive wording as requirements on programs or implementations.
  \item Explain when a violated requirement produces undefined behavior and when diagnostics remain required.
  \item Distinguish strictly conforming programs from the broader class of conforming programs.
  \item Compare hosted and freestanding implementation obligations.
  \item Evaluate whether an implementation extension preserves strict-program behavior.
  \item Identify the implementation-defined, locale-specific, and extension documentation that must accompany an implementation.
\end{itemize}

\section{Normative Structure Map}
\begin{CornellNotesTable}
\CornellNoteRow{Mandatory wording}{Requirements and prohibitions create obligations for programs and implementations.}
\CornellNoteRow{Strict conformance}{The portable subset avoids unspecified, undefined, and implementation-defined behavior.}
\CornellNoteRow{Hosted vs freestanding}{Hosted implementations accept full hosted programs; freestanding implementations support a defined minimum subset.}
\end{CornellNotesTable}

\section{Mandatory Wording and Violations}
\begin{CornellNotesTable}
\CornellNoteRow{What force does mandatory wording carry?}{A mandatory statement is a requirement on an implementation or program; its negative form is a prohibition.}
\CornellNoteRow{What happens when a requirement outside a constraint is violated?}{If the violated mandatory requirement is outside a constraint or runtime-constraint, the behavior is undefined.}
\CornellNoteRow{How else can undefined behavior be indicated?}{It can be stated explicitly or arise when the source intentionally provides no behavioral definition. These forms do not differ in emphasis.}
\CornellNoteRow{Can a construct be both a constraint violation and undefined?}{Yes. The undefined-behavior consequence does not cancel a diagnostic obligation imposed by the constraint rule.}
\CornellNoteRow{Can unspecified behavior appear in a correct program?}{Yes. A program can be correct while containing unspecified behavior, provided its executions obey the applicable observable-behavior rules.}
\end{CornellNotesTable}

\begin{CornellSummaryBox}{Section Summary}
Conformance analysis requires two parallel questions: what behavior follows, and whether a diagnostic is required. Undefined behavior does not erase an independently triggered diagnostic duty.
\end{CornellSummaryBox}



\section{Strictly Conforming and Conforming Programs}
\begin{CornellNotesTable}
\CornellNoteRow{What is a strictly conforming program?}{It uses only specified language and library features, stays within minimum implementation limits, and does not make its output depend on unspecified, undefined, or implementation-defined behavior.}
\CornellNoteRow{May a strict program use a conditional feature?}{Yes, when use is guarded by the associated feature test so that the program does not require an unavailable optional facility.}
\CornellNoteRow{What is a conforming program?}{A program accepted by some conforming implementation. It may depend on documented extensions or other nonportable features of that implementation.}
\CornellNoteRow{Why is strict conformance narrower?}{It targets maximum portability across conforming implementations rather than acceptability to one chosen implementation.}
\CornellNoteRow{What special rule applies to an active error directive?}{A preprocessing translation unit containing an active \texttt{\#error} directive must not be translated successfully; directives skipped by conditional inclusion do not trigger that prohibition.}
\end{CornellNotesTable}

\begin{CornellSummaryBox}{Section Summary}
Strict conformance is a portable subset; ordinary conformance is implementation-relative and can include documented nonportable dependencies.
\end{CornellSummaryBox}



\section{Hosted and Freestanding Implementations}
\begin{CornellNotesTable}
\CornellNoteRow{What must a hosted implementation accept?}{Every strictly conforming program intended for the hosted environment, together with the full set of required hosted library facilities.}
\CornellNoteRow{What must a freestanding implementation accept?}{Every strictly conforming program whose library use stays within the required freestanding subset, including specified foundational headers and the allowed portions of string and memory-alignment facilities.}
\CornellNoteRow{Is an operating system required for freestanding execution?}{No. Freestanding execution is designed to work without the services normally associated with a hosted operating environment.}
\CornellNoteRow{Can floating-point feature support expand the freestanding subset?}{Yes. When the associated floating-point feature macros are defined, specified environment, mathematics, and conversion facilities become usable under the stated restrictions.}
\CornellNoteRow{Do reserved identifiers change merely because execution is freestanding?}{Including a header carries the specified identifier-reservation consequences; the freestanding category does not generally turn reserved names into application names.}
\end{CornellNotesTable}

\begin{CornellSummaryBox}{Section Summary}
The categories differ mainly in required execution and library support. Freestanding does not mean unconstrained; it has a defined minimum acceptance contract.
\end{CornellSummaryBox}



\section{Extensions and Documentation}
\begin{CornellNotesTable}
\CornellNoteRow{Are extensions permitted?}{Yes. An implementation can add language or library facilities if doing so does not change the behavior of any strictly conforming program.}
\CornellNoteRow{What is the extension preservation test?}{Take a strictly conforming program and verify that the extension neither reinterprets its source nor changes its required behavior.}
\CornellNoteRow{What must implementation documentation define?}{All implementation-defined characteristics, locale-specific characteristics, and extensions supplied by that implementation.}
\CornellNoteRow{Why is documentation part of conformance?}{Programs can rely on implementation-defined behavior only relative to a documented implementation. Without the required documentation, that contract cannot be evaluated.}
\end{CornellNotesTable}

\begin{CornellSummaryBox}{Section Summary}
Extensions are compatible with conformance only when they preserve the specified meaning of strictly conforming programs and are documented as required.
\end{CornellSummaryBox}



\section{Programmer and Implementation Checklist}
\begin{itemize}
  \item[$\square$] Identify whether each mandatory statement applies to the program or implementation.
  \item[$\square$] Check whether a violation occurs inside a constraint, runtime-constraint, or neither.
  \item[$\square$] Preserve any diagnostic obligation even when behavior is also undefined.
  \item[$\square$] Reject strict-conformance claims that depend on unspecified, undefined, or implementation-defined output.
  \item[$\square$] Verify conditional features are protected by the corresponding feature tests.
  \item[$\square$] Classify the target implementation as hosted or freestanding before assessing required library support.
  \item[$\square$] For freestanding targets, compare every used header and function with the required subset.
  \item[$\square$] Test extensions against an unchanged strictly conforming program.
  \item[$\square$] Require documentation for implementation-defined and locale-specific choices and for extensions.
  \item[$\square$] Treat ordinary program conformance as relative to a particular conforming implementation.
\end{itemize}

\begin{CornellSummaryBox}{Clause Synthesis}
Clause 4 turns descriptive language into an accountability model. Strict programs avoid dependencies that would defeat cross-implementation portability. Hosted implementations accept the full strict subset; freestanding implementations accept a defined, smaller library subset. Extensions are allowed only if they preserve strict-program behavior, and implementation-specific choices must be documented. Diagnostics and behavioral consequences must be analyzed separately.
\end{CornellSummaryBox}



\clearpage
\section{Self-Test Questions}
\begin{enumerate}
  \item What does mandatory wording signify in this clause?
  \item When a mandatory rule outside a constraint or runtime-constraint is violated, what follows?
  \item Can an implementation translate a unit containing an active \texttt{\#error} directive successfully?
  \item May a strictly conforming program depend on documented implementation-defined output?
  \item How does a conforming program differ from a strictly conforming program?
  \item What are the two implementation categories?
  \item Does a freestanding implementation have to supply the full hosted library?
  \item A compiler adds a new keyword that makes an existing strict program parse differently. Is the extension automatically permitted?
  \item A constraint violation also has undefined behavior. Can the compiler omit the required diagnostic?
  \item What documentation must accompany a conforming implementation?
\end{enumerate}

\clearpage
\section{Answer Key}
\begin{enumerate}
  \item It imposes a requirement on a program or implementation; the negative form imposes a prohibition.

  \item The behavior is undefined under Clause 4.

  \item No. It must not do so. A directive in a group skipped by conditional inclusion is not active for this rule.

  \item No. Strict conformance excludes output dependent on implementation-defined behavior even when the implementation documents it.

  \item A conforming program need only be acceptable to a conforming implementation and can rely on that implementation's extensions or documented choices.

  \item Hosted and freestanding.

  \item No. It must supply the defined freestanding subset and any additional facilities triggered by supported conditional features.

  \item No. An extension must not alter the behavior of any strictly conforming program; changing the interpretation of existing strict source fails that preservation test.

  \item No. The constraint-based diagnostic obligation remains even when the behavior is also undefined.

  \item Documentation defining every implementation-defined and locale-specific characteristic and every extension.
\end{enumerate}

\section{Key-Term Recap}
\begin{CornellNotesTable}
\toprule
Term & Working meaning \\
\midrule
\endfirsthead
\toprule
Term & Working meaning (continued) \\
\midrule
\endhead
Strictly conforming program & A program using only specified features without output dependencies on unspecified, undefined, or implementation-defined behavior and without exceeding minimum limits. \\
Conforming program & A program acceptable to at least one conforming implementation. \\
Hosted implementation & An implementation providing the hosted execution model and required hosted facilities. \\
Freestanding implementation & An implementation supporting the freestanding model and its defined minimum library subset. \\
Extension & An added language or library facility that must preserve the behavior of strict programs. \\
Diagnostic obligation & A requirement to issue at least one diagnostic for an applicable syntax or constraint violation. \\
Conditional feature & An optional facility whose availability is indicated by an associated feature test. \\
\bottomrule
\end{CornellNotesTable}

\begin{CornellSummaryBox}{One-Sentence Takeaway}
Conformance is a three-way contract among program assumptions, implementation obligations, and documented choices, with strict conformance defining the maximally portable subset.
\end{CornellSummaryBox}
\end{document}
